\documentclass{article}
\usepackage[T1]{fontenc}
\usepackage{lmodern}
\usepackage{graphicx}
\usepackage{amsmath,amssymb}
\usepackage[a4paper,margin=2cm]{geometry}
\usepackage[absolute,overlay]{textpos}
\setlength{\TPHorizModule}{1mm}
\setlength{\TPVertModule}{1mm}
\usepackage[indonesian]{babel}
\usepackage{hyperref}
\setlength{\emergencystretch}{2em}
% unit: Halaman 12

\begin{document}

% === Halaman 12 (llm) ===

Contoh: Kurva eliptik $y^2 = x^3 + 2x + 4$

Misalkan $P(2, 4)$ dan $Q(0, 2)$ dua titik pada kurva

Penjumlahan titik: $P + Q = R$. Tentukan $R$!

Langkah-langkah menghitung koordinat $R$:

\begin{itemize}
    \item Gradien garis $g$: $m = (y_p - y_q)/(x_p - x_q) = (4 - 2)/(2 - 0) = 1$
    \item $x_r = m^2 - x_p - x_q = 1^2 - 2 - 0 = -1$
    \item $y_r = m(x_p - x_r) - y_p = 1(2 - (-1)) - 4 = -1$
    \item Jadi koordinat $R(-1, -1)$
    \item Periksa apakah $R(-1, -1)$ sebuah titik pada kurva eliptik:
    
    \begin{align*}
    y^2 = x^3 + 2x + 4 & \Leftrightarrow (-1)^2 = (-1)^3 + 2(-1) + 4 \\
    & \Leftrightarrow 1 = -1 - 2 + 4 \\
    & \Leftrightarrow 1 = 1 \text{ (terbukti } R(-1, -1) \text{ titik yang terletak pada kurva } y^2 = x^3 + 2x + 4)\
    \end{align*}
\end{itemize}

\end{document}
